\subsection{\texorpdfstring{应用：空间 \(\mathcal{M}_{+}^{b}(T)\) 的拓扑性质}{应用：空间 \textbackslash mathcal\{M\}\_\{+\}\^{}\{b\}(T) 的拓扑性质}}

首先注意，若 T 完全正则，则 \(\mathcal{M}^{b}(T)\) 是豪斯多夫拓扑向量空间，因而完全正则。因此，\(\mathcal{M}_{+}^{b}(T)\) 完全正则。

命题 10. --- 设 T 是波兰空间；则 \(\mathcal{M}_{+}^{b}(T)\) 在窄拓扑下是波兰空间。

先处理 T 是波兰空间且紧的情形。质量 \(\leqslant 1\) 的正测度所成的集合 U 是紧的（第 III 章 §1, No.~9, 命题 15 的推论 2），并且窄拓扑在 U 上诱导的拓扑（这里与淡拓扑一致）也由 \(\mathcal{C}(\mathrm{T})\) 的某个总子集上的逐点收敛拓扑诱导（同上，No.~10, 命题 17）。现在，\(\mathcal{C}(\mathrm{T})\) 中存在可数总集（GT, X, §3, No.~3, 定理 1）；因而 U 是可度量化的紧空间。质量 \textless{} 1 的正测度所成的集合 V 是 U 中的开集，故是波兰局部紧空间。现在，映射 \(\mu \mapsto \frac{1}{1 + \mu(1)}\mu\) 将 \(\mathcal{M}_{+}^{b}(\mathrm{T})\) 同胚映到 V，而映射 \(\lambda \mapsto \frac{1}{1 - \lambda(1)}\lambda\) 是其逆同胚。

现在处理 T 是波兰空间的情形；不妨设 T 是度量化紧空间 X 中递减开集序列 \((\mathbf{G}_{n})\) 的交（GT, IX, §6, No.~1, 定理 1 的推论 1）；于是 \(\mathcal{M}_{+}^{b}(\mathrm{T})\) 同胚于 \(\mathcal{M}_{+}^{b}(\mathrm{X})\) 的子空间 W，W 由集中于 T 的测度组成（第 3 号命题 8），只须证明 W 是波兰空间 \(\mathcal{M}_{+}^{b}(\mathrm{X})\) 中一列开集的交（GT, 同上，定理 1）。现在令 \(W_{n}\) 为集中于 \(\mu\in\mathcal{M}_{+}^{b}(\mathrm{X})\) 的测度 \(G_{n}\) 所成的集合；映射 \(h_{n}:\mu\mapsto\mu^{\bullet}(\mathrm{X}-\mathrm{G}_{n})\) 在 \(\mathcal{M}_{+}^{b}(\mathrm{X})\) 上上半连续（第 3 号命题 6），因而集合 \(\mathbf{A}_k^n\) 由满足 \(\mu \in \mathcal{M}_+^b(\mathbf{X})\) 的测度 \(h_n(\mu) < 1 / k\) 构成，对每个 \(k \geqslant 1\) 和每个 \(n \in \mathbf{N}\) 都是开集。注意到 \(\mathrm{W} = \bigcap_{n} \mathrm{W}_n = \bigcap_{n,k} \mathrm{A}_k^n\)，证明即完成。

推论 1. --- 若 T 是可数型的可度量化空间，则 \(\mathcal{M}_{+}^{b}(T)\) 在窄拓扑下是可度量化的可数型空间。

事实上，令 \(\widehat{T}\) 为赋予 T 拓扑的度量下 T 的完备化；空间 \(\widehat{T}\) 是波兰空间，而 \(\mathcal{M}_{+}^{b}(T)\) 同胚于波兰空间 \(\mathcal{M}_{+}^{b}(\widehat{T})\) 的子空间，该子空间由集中于 T 的测度组成（第 3 号命题 8）。但是，波兰空间的每个子空间都是可度量化的可数型空间（GT, IX, §2, No.~8）。

推论 2. --- 若 T 是完全正则的苏斯林空间（分别地，卢津空间），则 \(\mathcal{M}_{+}^{b}(\mathrm{T})\) 是苏斯林空间（分别地，卢津空间）。

事实上，取波兰空间 P 及从 P 到 T 的连续映射 f（GT, IX, §6, No.~2, 定义 2）。令 \(\widetilde{f}\) 为连续映射 \(\mu\mapsto f(\mu)\)，它从 \(\mathcal{M}_{+}^{b}(P)\) 到 \(\mathcal{M}_{+}^{b}(T)\)；由命题 10，空间 \(\mathcal{M}_{+}^{b}(P)\) 是波兰空间，而 \(\widetilde{f}\) 是满射（§2, No.~4, 命题 9）；所以 \(\mathcal{M}_{+}^{b}(T)\) 是苏斯林空间。同样，若 T 是卢津空间，则可假设 f 是单射（GT, 同上，No.~4, 命题 12）；于是 \(\widetilde{f}\) 是单射（§2, No.~4, 命题 8），从而 \(\mathcal{M}_{+}^{b}(T)\) 是卢津空间（GT, 同上，No.~4, 命题 12）。

设 T 是完全正则的苏斯林空间（回忆一下，只要 T 是苏斯林且正则的即可满足这一点（TG, 附录 1，命题 2 的推论）），并设 H 是 \(\mathcal{M}_{+}^{b}(T)\) 的紧子集；则 H 是紧的苏斯林空间，因而在窄拓扑下可度量化（同上，附录 1，命题 3 的推论 2）。

